\begin{abstract}
When humans encounter words like ``roses'' or ``thunder'' in text, they spontaneously activate sensory imagery---smelling the flowers, hearing the rumble---even without being asked to do so. Do large language models (LLMs) exhibit analogous behavior? We investigate whether LLMs spontaneously activate internal representations corresponding to sensory modalities when processing sensory-laden text, or whether such associations arise only post-hoc when explicitly prompted. Using \qwen, we extract hidden states for 150 words spanning six sensory modalities (auditory, gustatory, haptic, interoceptive, olfactory, visual) across three conditions: implicit narrative context, explicit sensory description, and neutral control. Linear probes achieve 90.0\% accuracy at decoding the dominant sensory modality from hidden states in the implicit condition---where no sensory language appears---demonstrating that LLMs spontaneously encode sensory information far above the 16.7\% chance level ($p = 0.005$). Implicit and explicit conditions activate overlapping representations (cosine similarity $= 0.822$, significantly higher than the control comparison, $p = 6.1 \times 10^{-20}$). The six senses form a structured geometry of mutually opposing directions in representation space, with interoception participating alongside the five classic senses. These findings demonstrate that text-only language models build a ``simulated sensory world'' from linguistic co-occurrence alone.
\end{abstract}
